% !TEX program = xelatex
% =====================================================================
%  BrewLite 1.0 — Báo cáo Assignment 1
%  Biên dịch: xelatex BaoCao_Assignment1.tex (chạy 2 lần để có mục lục)
% =====================================================================
\documentclass[12pt, a4paper]{article}

% ---------------------------------------------------------------------
% 1. LỀ TRANG
% ---------------------------------------------------------------------
\usepackage[a4paper, top=2.6cm, bottom=2.5cm, left=2.5cm, right=2.2cm,
            headheight=16pt, headsep=0.8cm, footskip=1.1cm]{geometry}

% ---------------------------------------------------------------------
% 2. NGÔN NGỮ & FONT CHỮ
%    Thân bài: Times New Roman (nếu máy có) -> TeX Gyre Termes (có sẵn trong TeX Live)
%    Tiêu đề:  TeX Gyre Heros (giống Helvetica/Arial)
%    Code:     DejaVu Sans Mono (hiển thị đủ dấu tiếng Việt)
% ---------------------------------------------------------------------
\usepackage{fontspec}
\usepackage[vietnamese, provide=*]{babel}
\IfFontExistsTF{Times New Roman}
  {\setmainfont{Times New Roman}}
  {\setmainfont{TeX Gyre Termes}}
\setsansfont{TeX Gyre Heros}
\setmonofont{DejaVu Sans Mono}[Scale=0.84]
\usepackage{microtype}
\usepackage{setspace}
\setstretch{1.18}

% ---------------------------------------------------------------------
% 3. GÓI CHỨC NĂNG
% ---------------------------------------------------------------------
\usepackage{amsmath, amssymb}
\usepackage[table]{xcolor}
\usepackage{array, tabularx, booktabs, colortbl}
\usepackage{enumitem}
\usepackage{titlesec}
\usepackage{fancyhdr}
\usepackage{caption}
\usepackage{float}
\usepackage{tikz}
\usetikzlibrary{babel}
\usepackage[most]{tcolorbox}
\usepackage{xurl}   % cho phép URL/đường dẫn dài tự xuống dòng
\usepackage[colorlinks=true, linkcolor=brandblue, urlcolor=brandblue,
            citecolor=brandblue, bookmarksnumbered=true,
            pdftitle={BrewLite 1.0 - Báo cáo Assignment 1},
            pdfauthor={Nhóm BrewLite}]{hyperref}

% ---------------------------------------------------------------------
% 4. BẢNG MÀU
% ---------------------------------------------------------------------
\definecolor{brandblue}{HTML}{1F3A5F}   % xanh navy chủ đạo
\definecolor{brandcoffee}{HTML}{8B5A2B} % nâu cà phê nhấn
\definecolor{brandcream}{HTML}{F7F1EA}  % nền kem nhạt
\definecolor{softgray}{HTML}{F3F5F8}
\definecolor{linegray}{HTML}{C9D1DC}
\definecolor{textgray}{HTML}{4A5565}
\definecolor{tagreq}{HTML}{B03A2E}      % BẮT BUỘC
\definecolor{tagplan}{HTML}{C27C0E}     % ĐỊNH HƯỚNG AS2
\definecolor{tagext}{HTML}{1E7F74}      % MỞ RỘNG

% ---------------------------------------------------------------------
% 5. ĐOẠN VĂN
% ---------------------------------------------------------------------
\setlength{\parindent}{0pt}
\setlength{\parskip}{6pt plus 1pt}
\setlength{\emergencystretch}{3em} % tránh chữ tràn lề / giãn xấu
\raggedbottom

% ---------------------------------------------------------------------
% 6. TIÊU ĐỀ MỤC
% ---------------------------------------------------------------------
\titleformat{\section}
  {\sffamily\Large\bfseries\color{brandblue}}
  {\colorbox{brandblue}{\makebox[1.6em]{\color{white}\thesection}}}{0.6em}{}
  [{\color{brandcoffee}\titlerule[1.2pt]}]
\titleformat{\subsection}
  {\sffamily\large\bfseries\color{brandblue}}{\thesubsection}{0.6em}{}
\titleformat{\subsubsection}
  {\sffamily\normalsize\bfseries\color{textgray}}{\thesubsubsection}{0.6em}{}
\titlespacing*{\section}{0pt}{18pt}{10pt}
\titlespacing*{\subsection}{0pt}{14pt}{6pt}
\titlespacing*{\subsubsection}{0pt}{10pt}{4pt}

% Mục lục
\setcounter{tocdepth}{3}

% ---------------------------------------------------------------------
% 7. DANH SÁCH
% ---------------------------------------------------------------------
\setlist{itemsep=2pt, topsep=4pt, parsep=0pt, leftmargin=1.6em}
\setlist[enumerate,1]{label=\textbf{\arabic*.}}
\setlist[itemize,1]{label={\color{brandcoffee}\small$\blacktriangleright$}}

% ---------------------------------------------------------------------
% 8. HEADER & FOOTER
% ---------------------------------------------------------------------
\pagestyle{fancy}
\fancyhf{}
\fancyhead[L]{\footnotesize\sffamily\color{textgray}\textit{BrewLite 1.0 — Báo Cáo Assignment 1: User Stories \& Cấu Trúc Dự Án}}
\fancyhead[R]{\footnotesize\sffamily\color{textgray}\textit{ĐH Sài Gòn — 2026}}
\fancyfoot[C]{\small\sffamily\color{textgray}\thepage}
\renewcommand{\headrule}{{\color{linegray}\hrule height 0.5pt}}
\renewcommand{\footrulewidth}{0pt}

% ---------------------------------------------------------------------
% 9. CHÚ THÍCH BẢNG
% ---------------------------------------------------------------------
\captionsetup{font={small,sf}, labelfont={bf,color=brandblue},
              skip=6pt, justification=centering}
\captionsetup[table]{position=below}

% ---------------------------------------------------------------------
% 10. CỘT BẢNG
% ---------------------------------------------------------------------
\newcolumntype{Y}{>{\raggedright\arraybackslash}X}
\newcolumntype{L}[1]{>{\raggedright\arraybackslash}p{#1}}
\newcolumntype{C}[1]{>{\centering\arraybackslash}p{#1}}
\renewcommand{\arraystretch}{1.35}

% ---------------------------------------------------------------------
% 11. MACRO DÙNG CHUNG
% ---------------------------------------------------------------------
% Nhãn phạm vi (badge)
\DeclareRobustCommand{\tagbox}[2]{%
  \tikz[baseline=(t.base)]\node[fill=#1, text=white, rounded corners=2pt,
     inner xsep=4pt, inner ysep=1.5pt, font=\sffamily\scriptsize\bfseries](t){#2};}
\newcommand{\tagReq}{\tagbox{tagreq}{BẮT BUỘC}}
\newcommand{\tagPlan}{\tagbox{tagplan}{ĐỊNH HƯỚNG AS2}}
\newcommand{\tagExt}{\tagbox{tagext}{MỞ RỘNG}}

% Đường dẫn file / mã trong văn bản (tự xuống dòng được)
\newcommand{\fpath}[1]{\nolinkurl{#1}}

% Thẻ User Story
%   #1 = mã (US01), #2 = tên, #3 = nhãn (\tagReq / \tagPlan / \tagExt)
\newtcolorbox{userstory}[3]{%
  enhanced, unbreakable,
  colback=white, colframe=linegray, boxrule=0.6pt, arc=3pt,
  left=10pt, right=10pt, top=6pt, bottom=6pt,
  before skip=10pt, after skip=10pt,
  colbacktitle=softgray, coltitle=brandblue,
  fonttitle=\sffamily\bfseries,
  borderline west={3pt}{0pt}{brandcoffee},
  title={#1.\enspace #2\hfill #3},
  toptitle=3pt, bottomtitle=3pt,
  titlerule=0.4pt, colupper=black,
  code={\phantomsection\addcontentsline{toc}{subsubsection}{#1. #2}}
}

% Các dòng bên trong thẻ
\newcommand{\usmeta}[2]{%
  {\small\sffamily\color{textgray}\textbf{Tác nhân:} #1\qquad\textbf{Phạm vi:} #2}\par\smallskip}
\newcommand{\usstory}[1]{%
  \begin{tcolorbox}[enhanced, colback=brandcream, colframe=brandcream,
      boxrule=0pt, arc=2pt, left=8pt, right=8pt, top=4pt, bottom=4pt,
      before skip=2pt, after skip=6pt]
    \textbf{User Story:} \textit{#1}
  \end{tcolorbox}}
\newcommand{\aclabel}{{\sffamily\bfseries\color{brandblue}Tiêu chí chấp nhận (AC):}\par}
\newcommand{\usimpl}[1]{%
  \par\smallskip{\color{linegray}\hrule height 0.4pt}\smallskip
  {\small\textbf{\color{brandcoffee}Hiện thực thực tế:} #1}\par}

% Hộp nhấn mạnh
\newtcolorbox{highlightbox}{%
  enhanced, colback=softgray, colframe=brandblue, boxrule=0pt,
  borderline west={3pt}{0pt}{brandblue}, arc=0pt,
  left=10pt, right=10pt, top=6pt, bottom=6pt, before skip=8pt, after skip=8pt}

% =====================================================================
\begin{document}
% =====================================================================

% ==================== TRANG BÌA ====================
\begin{titlepage}
\newgeometry{top=2cm, bottom=2cm, left=2.5cm, right=2.2cm}
\begin{tikzpicture}[remember picture, overlay]
  \draw[brandblue, line width=2pt]
    ([xshift=1.2cm, yshift=-1.2cm]current page.north west) rectangle
    ([xshift=-1.2cm, yshift=1.2cm]current page.south east);
  \draw[brandcoffee, line width=0.6pt]
    ([xshift=1.4cm, yshift=-1.4cm]current page.north west) rectangle
    ([xshift=-1.4cm, yshift=1.4cm]current page.south east);
\end{tikzpicture}

\begin{center}
  {\sffamily\bfseries\large\color{brandblue} TRƯỜNG ĐẠI HỌC SÀI GÒN}\\[3pt]
  {\sffamily\bfseries\large\color{brandblue} KHOA CÔNG NGHỆ THÔNG TIN}\\[3pt]
  {\sffamily\bfseries\normalsize\color{textgray} BỘ MÔN CÔNG NGHỆ PHẦN MỀM}\\[0.6cm]

  {\color{brandcoffee}\rule{0.35\textwidth}{1.2pt}}\\[1.6cm]

  {\sffamily\bfseries\Large\color{textgray} BÁO CÁO ASSIGNMENT 1}\\[0.5cm]
  {\sffamily\bfseries\fontsize{40}{46}\selectfont\color{brandblue} BREWLITE 1.0}\\[0.6cm]
  {\sffamily\bfseries\large\color{brandcoffee} USER STORIES, TIÊU CHÍ CHẤP NHẬN\\[2pt] VÀ CẤU TRÚC DỰ ÁN}\\[0.5cm]
  {\itshape\large Hệ thống đặt cà phê mộc và thanh toán không tiền mặt}\\[1.4cm]
\end{center}

\begin{tcolorbox}[enhanced, colback=softgray, colframe=softgray, arc=4pt,
                  left=12pt, right=12pt, top=10pt, bottom=10pt,
                  width=0.92\textwidth, center]
\small
\renewcommand{\arraystretch}{1.4}
\begin{tabularx}{\linewidth}{@{}L{4.4cm}Y@{}}
  \textbf{Học phần:}             & Công nghệ Phần mềm \\
  \textbf{Học kỳ:}               & Học kỳ 1 — Năm học 2026–2027 \\
  \textbf{Giảng viên hướng dẫn:} & Thầy Từ Lãng Phiêu \\
  \multicolumn{2}{@{}l@{}}{\color{linegray}\rule{\linewidth}{0.4pt}} \\[-4pt]
  \textbf{GitHub Repository:}    & \href{https://github.com/Longcodepro/Project_SE_Assignment01_HK1_2026_2027}{\fpath{Longcodepro/Project_SE_Assignment01_HK1_2026_2027}} \\
  \textbf{Frontend Live Web:}    & \url{https://se-as1-brewlite.vercel.app} \\
  \textbf{Backend Swagger UI:}   & \url{https://brewlite-backend.onrender.com/docs} \\
  \multicolumn{2}{@{}l@{}}{\color{linegray}\rule{\linewidth}{0.4pt}} \\[-4pt]
  \textbf{Nhóm thực hiện:}       & \textbf{Nhóm BrewLite (02 thành viên)} \\
  1.~\textbf{Nguyễn Thành Long}  & Nhóm trưởng — Backend, DevOps, Scrum Master \\
  2.~\textbf{Trần Tấn Tài}       & Frontend Lead, UI/UX Designer, Quản trị State \\
\end{tabularx}
\end{tcolorbox}

\vfill
\begin{center}
  {\sffamily\bfseries\color{brandblue} TP. HỒ CHÍ MINH — THÁNG 10 NĂM 2026}
\end{center}
\restoregeometry
\end{titlepage}

% ==================== MỤC LỤC ====================
\pagenumbering{roman}
\begingroup
\hypersetup{linkcolor=black}
\tableofcontents
\endgroup
\clearpage
\pagenumbering{arabic}

% ==================== NỘI DUNG CHÍNH ====================

\section{Thông Tin Và Phạm Vi Báo Cáo}

Báo cáo này trình bày toàn bộ nội dung của \textbf{Assignment 1} thuộc học phần Công nghệ Phần mềm: xây dựng hệ thống \textbf{User Stories} kèm tiêu chí chấp nhận (\textit{Acceptance Criteria}) chi tiết và xác lập \textbf{Cấu trúc dự án} chuẩn mực cho hệ thống bán lẻ cà phê không tiền mặt \textbf{BrewLite 1.0}.

Mã nguồn dự án được lưu trữ công khai tại GitHub: \href{https://github.com/Longcodepro/Project_SE_Assignment01_HK1_2026_2027}{\fpath{Longcodepro/Project_SE_Assignment01_HK1_2026_2027}}.

\subsection{Phạm vi nghiệp vụ cốt lõi}
Hệ thống giải quyết trọn vẹn luồng trải nghiệm khách hàng từ xem thực đơn, tùy biến kích cỡ (Size S/M/L) và topping, quản lý giỏ hàng, xác thực tài khoản JWT, tạo đơn hàng với giá do máy chủ tính độc quyền, thanh toán mô phỏng Thẻ/Ví, tra cứu lịch sử đơn hàng, cùng các quy tắc dữ liệu phân tán (chống gửi trùng Idempotency, khóa lạc quan version, tích điểm Loyalty, Order State Machine).

\subsection{Quy định giới hạn phạm vi theo chỉ đạo Giảng viên}
\begin{highlightbox}
Tuân thủ chỉ đạo của Giảng viên hướng dẫn, báo cáo Assignment 1 \textbf{tuyệt đối không đưa các sơ đồ phân tích thiết kế} như Sơ đồ Use Case, Activity Diagram, Sequence Diagram, Class Diagram hay tài liệu kiến trúc phần mềm vào báo cáo này (toàn bộ các sơ đồ trên thuộc phạm vi yêu cầu của \textbf{Assignment 2}).
\end{highlightbox}

\subsection{Tuyên ngôn thực thi của nhóm}
Khác với các đề xuất lý thuyết trên giấy, \textbf{100\% User Stories trong báo cáo này đã được nhóm hiện thực hóa thành mã nguồn hoạt động thật}, vượt qua 52 Pull Requests tự động hóa CI/CD trên GitHub Actions và đã triển khai vận hành trực tuyến 24/7 trên môi trường Đám mây.

% =====================================================================
\section{User Stories Và Tiêu Chí Chấp Nhận (Acceptance Criteria)}

Mỗi User Story được viết theo cấu trúc chuẩn: \textit{“Là [tác nhân], tôi muốn [chức năng], để [giá trị]”} kèm các tiêu chí chấp nhận (\textit{Acceptance Criteria -- AC}) định lượng, kiểm chứng được.

\subsection[Nhóm User Stories Cốt Lõi {[}BẮT BUỘC{]}]{Nhóm User Stories Cốt Lõi \texorpdfstring{\tagReq}{[BẮT BUỘC]}}

% --- US01 ---
\begin{userstory}{US01}{Xem thực đơn đồ uống}{\tagReq}
\usmeta{Khách hàng.}{Bắt buộc.}
\usstory{Là khách hàng, tôi muốn xem danh sách các món đồ uống hiển thị kèm hình ảnh, tên món và giá niêm yết để lựa chọn sản phẩm phù hợp.}
\aclabel
\begin{enumerate}
  \item API trả danh sách sản phẩm gồm mã (\texttt{id}), tên (\texttt{name}), giá chuẩn (\texttt{price}) và ảnh (\texttt{imageUrl}).
  \item Giao diện thể hiện đủ 3 trạng thái: Đang tải (\textit{Loading}), Danh sách rỗng (\textit{Empty}) và Lỗi kết nối (\textit{Error}).
  \item Cung cấp nút thử lại (\textit{Retry}) hoặc công cụ cấu hình kết nối khi máy chủ phản hồi chậm hoặc ngắt kết nối.
\end{enumerate}
\usimpl{API \texttt{GET /products}, trang \fpath{frontend/src/app/page.tsx}, test \fpath{products.controller.spec.ts}.}
\end{userstory}

% --- US02 ---
\begin{userstory}{US02}{Xem chi tiết và tùy chỉnh đồ uống}{\tagReq}
\usmeta{Khách hàng.}{Bắt buộc.}
\usstory{Là khách hàng, tôi muốn mở hộp thoại chi tiết để chọn kích cỡ (Size) và các loại topping thêm vào, giúp đồ uống hợp khẩu vị cá nhân.}
\aclabel
\begin{enumerate}
  \item Chỉ hiển thị kích cỡ hợp lệ (S: +0đ, M: +5.000đ, L: +10.000đ) và topping hợp lệ (Trân châu: 8.000đ, Kem cheese: 10.000đ, Thạch cafe: 6.000đ).
  \item Giá tạm tính trên giao diện tự động cập nhật chính xác theo thời gian thực mỗi khi thay đổi tùy chọn hoặc số lượng.
  \item Chặn thêm vào giỏ nếu số lượng không nằm trong khoảng $1 \le \text{quantity} \le 99$.
\end{enumerate}
\usimpl{Component \fpath{ProductDetailModal.tsx}, quy tắc bảng giá \fpath{backend/src/orders/pricing.ts}.}
\end{userstory}

% --- US03 ---
\begin{userstory}{US03}{Quản lý giỏ hàng}{\tagReq}
\usmeta{Khách hàng.}{Bắt buộc.}
\usstory{Là khách hàng, tôi muốn quản lý giỏ hàng (thay đổi số lượng, xóa món) để kiểm tra tổng chi phí trước khi đặt hàng.}
\aclabel
\begin{enumerate}
  \item Cho phép thêm, tăng/giảm số lượng và xóa từng dòng món trong giỏ.
  \item Hai món cùng loại nhưng khác cấu hình Size hoặc Topping được phân biệt thành hai dòng riêng biệt.
  \item Tổng số lượng trên huy hiệu giỏ hàng và tổng tiền cập nhật tức thì; giỏ rỗng thì nút thanh toán bị vô hiệu hóa.
  \item Dữ liệu giỏ hàng lưu trữ bền vững qua \texttt{localStorage} để không bị mất khi tải lại trang.
\end{enumerate}
\usimpl{Store Zustand \fpath{frontend/src/store/useCartStore.ts}, trang \fpath{frontend/src/app/cart/page.tsx}.}
\end{userstory}

% --- US04 ---
\begin{userstory}{US04}{Đăng ký và đăng nhập bảo mật}{\tagReq}
\usmeta{Khách hàng.}{Bắt buộc.}
\usstory{Là khách hàng, tôi muốn đăng ký và đăng nhập an toàn để thông tin cá nhân và đơn hàng được bảo vệ.}
\aclabel
\begin{enumerate}
  \item Đăng ký yêu cầu email hợp lệ và mật khẩu $\ge 6$ ký tự; từ chối đăng ký trùng email (\texttt{409 Conflict}).
  \item Mật khẩu băm một chiều bằng \texttt{bcryptjs} ($\mathit{salt} = 10$), tuyệt đối không lưu mật khẩu thuần.
  \item Đăng nhập thành công trả về JWT Bearer Token; các API nhạy cảm từ chối nếu thiếu token (\texttt{401 Unauthorized}).
\end{enumerate}
\usimpl{Module \fpath{backend/src/auth/}, \texttt{JwtAuthGuard}, modal \fpath{AuthModal.tsx}, 9 unit tests.}
\end{userstory}

% --- US05 ---
\begin{userstory}{US05}{Tạo đơn hàng với giá tính tại Máy chủ}{\tagReq}
\usmeta{Khách hàng đã đăng nhập.}{Bắt buộc.}
\usstory{Là khách hàng đã đăng nhập, tôi muốn gửi yêu cầu đặt hàng từ giỏ để hệ thống ghi nhận đơn vào cơ sở dữ liệu.}
\aclabel
\begin{enumerate}
  \item Client chỉ gửi mã món, số lượng, size và topping; \textbf{máy chủ tự tính toán lại toàn bộ tổng tiền}.
  \item Mọi tham số lạ hoặc can thiệp giá từ Client đều bị chặn bởi \texttt{class-validator} (\texttt{400 Bad Request}).
  \item Đơn hàng được tạo với trạng thái khởi tạo duy nhất là \texttt{PENDING}, sinh mã đơn duy nhất và gắn với người dùng.
\end{enumerate}
\usimpl{API \texttt{POST /orders}, DTO validation \fpath{create-order.dto.ts}, 15 unit tests.}
\end{userstory}

% --- US06 ---
\begin{userstory}{US06}{Thanh toán không tiền mặt bằng Ví hoặc Thẻ}{\tagReq}
\usmeta{Khách hàng.}{Bắt buộc.}
\usstory{Là khách hàng, tôi muốn thanh toán bằng Thẻ ngân hàng hoặc Ví điện tử để hoàn tất giao dịch không tiền mặt.}
\aclabel
\begin{enumerate}
  \item Khách hàng chọn phương thức thanh toán: Ví điện tử (\texttt{WALLET}) hoặc Thẻ ngân hàng (\texttt{CARD}).
  \item Quy tắc mô phỏng: Thẻ có 4 số cuối là \texttt{0000} bị từ chối; Ví điện tử từ chối khi đơn hàng vượt quá 500.000 VNĐ.
  \item Thanh toán thành công chuyển đơn sang \texttt{PAID}; thất bại chuyển sang \texttt{PAYMENT\_FAILED} (được thanh toán lại).
\end{enumerate}
\usimpl{API \texttt{POST /payments}, dịch vụ \fpath{payments.service.ts}, 24 unit tests.}
\end{userstory}

% --- US07 ---
\begin{userstory}{US07}{Xem màn hình xác nhận đơn hàng}{\tagReq}
\usmeta{Khách hàng.}{Bắt buộc.}
\usstory{Là khách hàng, tôi muốn nhận màn hình xác nhận chi tiết sau khi thanh toán để biết đơn đã được ghi nhận.}
\aclabel
\begin{enumerate}
  \item Hiển thị rõ mã đơn hàng, danh sách món, phụ thu, phương thức thanh toán và tổng số tiền thanh toán thực tế.
  \item Tải lại trang không tạo đơn hàng mới hay thanh toán lại; chỉ hiển thị trạng thái đã được máy chủ xác nhận.
\end{enumerate}
\usimpl{Màn hình \fpath{frontend/src/app/checkout/page.tsx}, cơ chế lưu vết đơn hàng.}
\end{userstory}

% --- US08 ---
\begin{userstory}{US08}{Xem lịch sử đơn hàng cá nhân}{\tagReq}
\usmeta{Khách hàng.}{Bắt buộc.}
\usstory{Là khách hàng, tôi muốn xem lại danh sách các đơn hàng đã đặt để theo dõi trạng thái và lịch sử chi tiêu.}
\aclabel
\begin{enumerate}
  \item Khách hàng chỉ xem được danh sách đơn hàng gắn liền với tài khoản của chính mình (bảo vệ bằng JWT claim).
  \item Danh sách hiển thị theo thứ tự thời gian mới nhất, kèm mã đơn, thời gian tạo, tổng tiền và nhãn trạng thái (\texttt{PENDING}, \texttt{PAID}, \texttt{PAYMENT\_FAILED}).
\end{enumerate}
\usimpl{API \texttt{GET /orders/me}, trang giao diện \fpath{frontend/src/app/orders/page.tsx}.}
\end{userstory}

% --- US09 ---
\begin{userstory}{US09}{Chống giao dịch trùng lặp (Distributed Idempotency)}{\tagReq}
\usmeta{Hệ thống.}{Bắt buộc.}
\usstory{Là hệ thống, tôi muốn nhận diện và xử lý yêu cầu gửi trùng lặp để khách hàng không bị trừ tiền hai lần.}
\aclabel
\begin{enumerate}
  \item Client gửi kèm Header \texttt{Idempotency-Key} trong yêu cầu tạo đơn (\texttt{POST /orders}) và thanh toán (\texttt{POST /payments}).
  \item Yêu cầu lặp lại có cùng khóa \texttt{Idempotency-Key} và cùng tài khoản được trả về kết quả cũ mà không thanh toán lại.
  \item Cột \texttt{idempotencyKey} được đánh chỉ mục và ràng buộc duy nhất (\texttt{@@unique([userId, idempotencyKey])}).
\end{enumerate}
\usimpl{Migration \fpath{them_idempotency_key}, xử lý tại \fpath{orders.service.ts} và \fpath{payments.service.ts}.}
\end{userstory}

% --- US10 ---
\begin{userstory}{US10}{Khóa lạc quan chống tranh chấp dữ liệu (Optimistic Locking)}{\tagReq}
\usmeta{Hệ thống.}{Bắt buộc.}
\usstory{Là hệ thống, tôi muốn kiểm soát các phiên cập nhật đồng thời lên cùng một đơn hàng để dữ liệu không bị ghi đè.}
\aclabel
\begin{enumerate}
  \item Bảng \texttt{Order} chứa cột \texttt{version} kiểu số nguyên, tự động tăng 1 mỗi khi đơn hàng thay đổi trạng thái.
  \item Mọi câu lệnh cập nhật bắt buộc phải khớp chính xác phiên bản \texttt{version} đã đọc trước đó (\texttt{where: \{ id, version \}}).
  \item Nếu phiên bản bị sai lệch do tiến trình khác cập nhật trước, hệ thống từ chối cập nhật và trả về mã lỗi \texttt{409 Conflict}.
\end{enumerate}
\usimpl{Migration \fpath{them_version_khoa_lac_quan}, logic cập nhật tại \fpath{payments.service.ts}.}
\end{userstory}

% --- US11 ---
\begin{userstory}{US11}{Quản lý khuyến mãi và chiết khấu}{\tagPlan}
\usmeta{Khách hàng.}{Định hướng Assignment 2.}
\usstory{Là khách hàng, tôi muốn áp dụng mã giảm giá hợp lệ để được giảm trừ trực tiếp trên tổng giá trị thanh toán.}
\aclabel
\begin{enumerate}
  \item Máy chủ kiểm tra điều kiện áp dụng mã (hạn dùng, giá trị đơn tối thiểu); mã không hợp lệ trả về thông báo rõ ràng.
  \item Số tiền giảm được khấu trừ nhất quán ở bước thanh toán và lưu vết vào chi tiết hóa đơn.
\end{enumerate}
\end{userstory}

% --- US12 ---
\begin{userstory}{US12}{Tích lũy điểm thưởng thành viên (Loyalty Engine)}{\tagReq}
\usmeta{Khách hàng.}{Bắt buộc.}
\usstory{Là khách hàng, tôi muốn được tự động tích lũy điểm thưởng sau mỗi lần thanh toán thành công để nhận quyền lợi.}
\aclabel
\begin{enumerate}
  \item Điểm thưởng chỉ được cộng vào tài khoản người dùng sau khi đơn hàng chuyển sang trạng thái \texttt{PAID} thành công.
  \item Tỷ lệ quy đổi chuẩn: Tích lũy 1 điểm cho mỗi 1.000 VNĐ chi tiêu thành công (ví dụ đơn 52.000đ tích được 52 điểm).
  \item Một đơn hàng không bao giờ được cộng điểm nhiều lần; điểm lưu trong cột \texttt{loyaltyPoints} của bảng \texttt{User}.
\end{enumerate}
\usimpl{Migration \fpath{them_tich_diem}, module \fpath{backend/src/payments/loyalty.ts}, 6 unit tests.}
\end{userstory}

% --- US13 ---
\begin{userstory}{US13}{Máy trạng thái bảo đảm vòng đời đơn hợp lệ (State Machine)}{\tagReq}
\usmeta{Hệ thống.}{Bắt buộc.}
\usstory{Là hệ thống, tôi muốn kiểm soát chặt chẽ việc chuyển đổi trạng thái của đơn hàng nhằm bảo đảm quy trình vận hành.}
\aclabel
\begin{enumerate}
  \item Ma trận trạng thái chỉ chấp nhận: \texttt{PENDING} $\rightarrow$ \texttt{PAID} hoặc \texttt{PAYMENT\_FAILED}; \texttt{PAYMENT\_FAILED} $\rightarrow$ \texttt{PAID} (thanh toán lại).
  \item Đơn hàng đã ở trạng thái \texttt{PAID} là trạng thái kết thúc, tuyệt đối cấm chuyển tiếp; vi phạm trả về \texttt{400 Bad Request}.
\end{enumerate}
\usimpl{Bộ điều phối máy trạng thái \fpath{backend/src/orders/order-status.ts}, 7 unit tests.}
\end{userstory}

\subsection[Nhóm User Stories Mở Rộng Dành Cho Barista {[}MỞ RỘNG{]}]{Nhóm User Stories Mở Rộng Dành Cho Barista \texorpdfstring{\tagExt}{[MỞ RỘNG]}}

\begin{tcolorbox}[enhanced, breakable, colback=white, colframe=linegray, boxrule=0.6pt,
    arc=3pt, borderline west={3pt}{0pt}{tagext}, left=10pt, right=10pt, top=6pt, bottom=6pt]
\begin{itemize}[itemsep=5pt]
  \item \textbf{US14. Tiếp nhận đơn chờ pha} \tagExt: Barista xem danh sách các đơn đã \texttt{PAID} và bấm tiếp nhận pha chế (\texttt{PREPARING}).
  \item \textbf{US15. Hoàn tất pha chế và giao đồ uống} \tagExt: Barista cập nhật trạng thái đơn thành \texttt{READY} (đã pha) và \texttt{COMPLETED} (đã giao).
  \item \textbf{US16. Báo sự cố và hủy đơn} \tagExt: Barista hủy đơn khi hết nguyên liệu đột xuất kèm lý do hủy bắt buộc.
\end{itemize}
\end{tcolorbox}

% =====================================================================
\section{Cấu Trúc Thư Mục Dự Án (Repository Structure)}

Cấu trúc dưới đây mô tả chính xác 100\% tổ chức mã nguồn thực tế của kho lưu trữ GitHub \fpath{Longcodepro/Project_SE_Assignment01_HK1_2026_2027}:

% Cây thư mục: cột trái là đường dẫn, cột phải là mô tả (mô tả dài tự xuống dòng thẳng hàng)
\newlength{\dirw}\setlength{\dirw}{7.5cm}
\newcommand{\dirrow}[2]{\par\noindent\hangindent=\dirw\hangafter=1
  \makebox[\dirw][l]{\footnotesize\ttfamily #1}{\small\color{textgray}#2}\par}
\newcommand{\dirgap}{\par\noindent{\footnotesize\ttfamily |}\par}

\begin{tcolorbox}[enhanced, breakable, colback=softgray, colframe=linegray, boxrule=0.5pt,
    arc=3pt, left=10pt, right=10pt, top=6pt, bottom=6pt,
    title={\sffamily\bfseries\small Project\_SE\_Assignment01\_HK1\_2026\_2027/},
    colbacktitle=brandblue, coltitle=white]
\raggedright\setlength{\parskip}{0pt}\setstretch{1.05}
\dirrow{|-- frontend/}{Ứng dụng Client (Next.js 16 App Router)}
\dirrow{|\ \ \ |-- src/app/}{Các trang: Menu, Giỏ hàng, Thanh toán, Đơn hàng}
\dirrow{|\ \ \ |-- src/components/}{UI dùng chung: Navbar, AuthModal, ProductDetailModal}
\dirrow{|\ \ \ |-- src/services/api.ts}{Axios Client, dynamic URL, token injection}
\dirrow{|\ \ \ |-- src/store/}{Quản lý State bằng Zustand (useCartStore, useAuthStore)}
\dirrow{|\ \ \ |-- src/types/index.ts}{Định nghĩa kiểu dữ liệu TypeScript dùng chung}
\dirrow{|\ \ \ `-- Dockerfile}{Đóng gói Frontend container (Node 22)}
\dirgap
\dirrow{|-- backend/}{Máy chủ API (NestJS 12 + Prisma 7)}
\dirrow{|\ \ \ |-- src/auth/}{Module Xác thực JWT \& băm bcrypt (US04)}
\dirrow{|\ \ \ |-- src/orders/}{Module Tạo đơn, State Machine, Khóa lạc quan}
\dirrow{|\ \ \ |-- src/payments/}{Module Thanh toán Ví/Thẻ \& Tích điểm}
\dirrow{|\ \ \ |-- src/products/}{Module Danh mục sản phẩm \& hình ảnh (US01)}
\dirrow{|\ \ \ |-- src/prisma/}{Module kết nối cơ sở dữ liệu toàn cục}
\dirrow{|\ \ \ |-- src/main.ts}{Điểm khởi chạy, HTTP Logger, CORS đa nguồn}
\dirrow{|\ \ \ |-- prisma/schema.prisma}{Lược đồ quan hệ thực thể cơ sở dữ liệu}
\dirrow{|\ \ \ |-- prisma/migrations/}{7 bản migration cơ sở dữ liệu có sẵn}
\dirrow{|\ \ \ `-- Dockerfile}{Đóng gói Backend container (Node 22)}
\dirgap
\dirrow{|-- docs/}{Hệ thống tài liệu dự án tinh gọn}
\dirrow{|\ \ \ |-- HUONG-DAN-TRIEN-KHAI-VA-CHAY.md}{Cẩm nang vận hành 4 kịch bản \& kinh nghiệm}
\dirrow{|\ \ \ |-- NHAT-KY-COMMIT.md}{Toàn bộ lịch sử commit tiếng Việt có dấu}
\dirrow{|\ \ \ |-- NHAT-KY-TIEN-DO.md}{Bảng tổng kết 52 PRs \& Retrospective}
\dirrow{|\ \ \ |-- TienDoCongViec.pdf}{Báo cáo tiến độ học thuật (XeLaTeX 14 trang)}
\dirrow{|\ \ \ `-- BaoCao\_Assignment1.tex}{Mã nguồn XeLaTeX báo cáo Assignment 1}
\dirgap
\dirrow{|-- scripts/}{Kịch bản tự động hóa hỗ trợ vận hành}
\dirrow{|\ \ \ |-- logs.sh}{Xem log HTTP request thời gian thực từ Docker}
\dirrow{|\ \ \ |-- start-tunnel.sh}{Chạy Cloudflare Tunnel qua giao thức HTTP/2}
\dirrow{|\ \ \ `-- studio-cloud.sh}{Mở Prisma Studio kết nối Database Cloud Neon}
\dirgap
\dirrow{|-- .github/workflows/}{CI/CD Workflows tự động kiểm tra mã nguồn}
\dirrow{|-- docker-compose.yml}{Thiết lập cụm dịch vụ cục bộ (Cổng 4000, 4001, 4002)}
\dirrow{`-- README.md}{Cẩm nang toàn tập hướng dẫn dự án chuẩn quốc tế}
\end{tcolorbox}

% =====================================================================
\clearpage
\section{Phân Công Nội Dung Và Trách Nhiệm Thành Viên}

Dự án được thực hiện bởi nhóm 2 thành viên theo mô hình \textbf{Agile Scrum}, phân chia công việc độc lập nhưng phối hợp chặt chẽ:

\begin{table}[H]
\centering
\small
\rowcolors{2}{softgray}{white}
\begin{tabularx}{\textwidth}{@{}L{3.1cm} Y L{4.0cm} C{1.2cm}@{}}
\rowcolor{brandblue}
\color{white}\textbf{Thành viên} & \color{white}\textbf{Nội dung \& User Story phụ trách} & \color{white}\textbf{Khu vực mã nguồn liên quan} & \color{white}\textbf{Tỷ lệ} \\
\midrule
\textbf{Nguyễn Thành Long}\newline{\footnotesize\color{textgray}(Nhóm trưởng)} &
Khởi tạo Monorepo, Docker Compose, CI/CD GitHub Actions. Hiện thực US01 (API), US04 (Backend Auth), US05 (API Orders), US06 (API Payments), US09 (Idempotency), US10 (Khóa lạc quan), US12 (Tích điểm), US13 (State Machine). Triển khai Render \& Neon Cloud. Soạn thảo XeLaTeX. &
\fpath{backend/src/}\newline \fpath{backend/prisma/}\newline \fpath{.github/workflows/}\newline \fpath{docker-compose.yml}\newline \fpath{scripts/}\newline \fpath{docs/} & \textbf{50\%} \\
\textbf{Trần Tấn Tài} &
Khởi tạo Next.js 16 (App Router), Design System Tailwind CSS. Hiện thực US01 (UI Menu), US02 (Modal chi tiết món), US03 (Giỏ hàng Zustand/LocalStorage), US04 (UI AuthModal), US07 (UI Xác nhận đơn), US08 (UI Lịch sử đơn). Tối ưu Mobile Responsive, Dark Mode, Server Config Modal \& Ping test. &
\fpath{frontend/src/app/}\newline \fpath{frontend/src/components/}\newline \fpath{frontend/src/store/}\newline \fpath{frontend/src/services/} & \textbf{50\%} \\
\bottomrule
\end{tabularx}
\caption{Bảng phân công trách nhiệm thực hiện Assignment 1}
\end{table}

\subsection{Quy trình phối hợp và kiểm soát chất lượng (Quality Assurance)}
\begin{enumerate}[itemsep=4pt]
  \item \textbf{Branching Strategy:} Mỗi tính năng là một nhánh riêng mang tiền tố tên người làm (\texttt{long/<task>} hoặc \texttt{tai/<task>}). Không bao giờ commit trực tiếp lên nhánh \texttt{main}.
  \item \textbf{Conventional Commits:} Toàn bộ commit tuân thủ quy ước chuẩn bằng tiếng Việt có dấu (ví dụ: \texttt{feat(task-6): thêm API POST /orders...}).
  \item \textbf{Pull Request \& Review chéo:} Mỗi nhánh khi hoàn thành đều mở PR, bắt buộc phải vượt qua 100\% các bài test tự động của CI Actions và được thành viên còn lại duyệt mới được merge.
\end{enumerate}

% =====================================================================
\clearpage
\section{Các Điểm Sáng \& Ưu Thế Vượt Trội Của Đội Ngũ (Showcase)}

So với các tiêu chuẩn thông thường của Assignment 1, nhóm BrewLite đã đạt được những kết quả vượt trội:

\begin{table}[H]
\centering
\small
\rowcolors{2}{softgray}{white}
\begin{tabularx}{\textwidth}{@{}L{3.3cm} Y Y@{}}
\rowcolor{brandblue}
\color{white}\textbf{Tiêu chí so sánh} & \color{white}\textbf{Mức độ thông thường} & \color{white}\textbf{Kết quả thực tế của Nhóm BrewLite} \\
\midrule
\textbf{1. Hoàn thiện mã nguồn} &
Chỉ đề xuất cấu trúc thư mục trên văn bản Word; mã nguồn chưa chạy hoặc chạy dở dang. &
\textbf{100\% mã nguồn đã chạy thật}, hoàn thành đầy đủ cả 3 Sprint, tích hợp toàn diện Frontend và Backend. \\
\textbf{2. Kiểm thử tự động (CI/CD)} &
Kiểm thử thủ công bằng tay; không có quy trình tự động hóa. &
Tự động hóa hoàn toàn với \textbf{52 Pull Requests đã merge}, 100\% vượt qua 64 unit tests (Vitest), kiểm tra cú pháp (Oxlint, ESLint) và Docker build. \\
\textbf{3. Xử lý phân tán chuyên sâu} &
Chỉ có CRUD cơ bản, dễ bị lỗi trừ tiền 2 lần hoặc xung đột ghi dữ liệu. &
Hiện thực đầy đủ các chuẩn công nghiệp cao cấp: \textbf{Idempotency-Key}, \textbf{Khóa lạc quan (version)}, \textbf{State Machine} và \textbf{Tính giá Server độc quyền}. \\
\textbf{4. Vận hành đa môi trường} &
Chỉ chạy được trên máy cục bộ của người lập trình; phụ thuộc vào cấu hình máy. &
Cung cấp \textbf{4 kịch bản vận hành}: Docker Compose 1-chạm, Local Dev, Hybrid Tunnel cho điện thoại và \textbf{100\% Cloud Production 24/7} (Vercel + Render + Neon). \\
\textbf{5. Đổi Server API động} &
URL Backend bị nhúng cứng; khó test khi mang sang máy khác hoặc trên điện thoại. &
Tích hợp nút \textbf{Server (cột sóng)} ngay trên Navbar web, lưu \texttt{localStorage}, cho phép đổi link API và \textbf{Ping test độ trễ trực tiếp}. \\
\textbf{6. Định dạng tài liệu học thuật} &
Soạn thảo trên Microsoft Word định dạng đơn giản. &
Soạn thảo và biên dịch bằng \textbf{XeLaTeX chuyên nghiệp}, font chữ sắc nét, cấu trúc phân cấp học thuật chuẩn chỉ. \\
\bottomrule
\end{tabularx}
\caption{Bảng đối chiếu ưu thế công nghệ vượt trội của nhóm BrewLite}
\end{table}

% =====================================================================
\section{Kết Luận Và Hướng Phát Triển Assignment 2}

Báo cáo Assignment 1 của nhóm BrewLite đã hoàn thành trọn vẹn và vượt mức toàn bộ các yêu cầu đặt ra:
\begin{itemize}[itemsep=4pt]
  \item Hệ thống 13 User Stories cốt lõi được mô tả chuẩn mực, định lượng và có tiêu chí chấp nhận kiểm chứng được.
  \item Cấu trúc thư mục Monorepo trong sáng, module hóa cao và sẵn sàng mở rộng.
  \item Nền tảng kỹ thuật vững chắc với 52 PRs và hệ thống triển khai đám mây thực tế 24/7.
\end{itemize}

\begin{highlightbox}
Đây là tiền đề kỹ thuật hoàn hảo để nhóm tự tin bước vào \textbf{Assignment 2}: xây dựng tài liệu đặc tả yêu cầu phần mềm (SRS), mô hình hóa hệ thống bằng hệ thống sơ đồ PlantUML (Use Case, Activity, Sequence, Class Diagram), thiết kế tài liệu kiến trúc (ADR) và mở rộng các tính năng phân hệ quản lý Barista.
\end{highlightbox}

\end{document}